\section{Language como interfaz cognitiva unificada}

\subsection{Por qué Language sigue siendo el medio central}
El curso señala que el lenguaje tiene un triple papel: medio de codificación de estado, medio de expresión de políticas, medio de protocolo de colaboración. Esto permite que el lenguaje no solo transmita explicaciones en lenguaje natural, sino que también sea vehículo de estructuras simbólicas, fragmentos de programa y registros de decisiones intermedias.

\begin{figure}[H]
\centering
\includegraphics[width=0.84\textwidth]{slide-03-reasoning-taxonomy.jpg}
\caption{Diapositiva del video, aproximadamente en 00:37:20: el curso organiza por capas los tipos de razonamiento según su punto de aplicación y su objetivo}
\end{figure}

\begin{importantbox}{Ventajas sistémicas de la interfaz de lenguaje}
\begin{itemize}
    \item Unificación entre modalidades: el texto puede servir como capa de intercambio entre visión, código y datos estructurados;
    \item Unificación entre módulos: el planner, el memory manager y el tool router pueden compartir el mismo protocolo;
    \item Unificación entre agentes: la colaboración multi-agent puede prototiparse primero mediante un protocolo de lenguaje.
\end{itemize}
\end{importantbox}

\subsection{La fusión entre el razonamiento simbólico y el razonamiento en lenguaje}
Yu Su no contrapone symbolic reasoning y neural reasoning, sino que enfatiza una «división del trabajo fusionada»: el razonamiento en lenguaje se encarga de la compresión semántica de alta dimensión y de la exploración de políticas, mientras que el razonamiento simbólico se encarga de las restricciones precisas y del cálculo verificable. La combinación de ambos es determinante para la estabilidad de las tareas de cadena larga.

\begin{knowledgebox}{Ejemplos de estrategias de fusión}
\begin{itemize}
    \item El modelo de lenguaje genera planes candidatos y el módulo simbólico realiza la verificación de restricciones;
    \item El modelo de lenguaje realiza búsqueda heurística y el ejecutor de programas hace la verificación final;
    \item El modelo de lenguaje explica los errores y el sistema de reglas hace la intercepción de seguridad.
\end{itemize}
\end{knowledgebox}

\begin{warningbox}{Los límites de depender solo del razonamiento en lenguaje}
Cuando la tarea requiere alta precisión numérica, optimización combinatoria con restricciones fuertes o demostración formal, la cadena de lenguaje natural pura tiende a distorsionarse, y es necesario introducir una representación intermedia programática.
\end{warningbox}

\subsection{Resumen de la sección}
Language no es un «formato de entrada de prompts», sino la capa de comunicación unificada del sistema de Agent. Su valor está en conectar los distintos módulos de capacidad, no en sustituir a todos los módulos.

